% ===================================================================
%  तक्ता क्रमांक ६ — घराचा आतला भाग, सामानसुमान, अन्न, प्रकाश.
%
%  Traduction, faite en 2026, en marathi d'aujourd'hui, sur l'ido de
%  texto/io. Regles de la colonne : voir 10-takta-01.tex. Cinq
%  scenes — la chambre, le salon, la salle a manger, la cuisine, la
%  salle de bain — et l'auteur remet la numerotation a 1 a chaque
%  fois : les cles portent c1 a c5.
%
%  DEUX NOTES SUR LA MEME PAGE, DEUX \VUnotes. Le fac-simile les
%  compose dans un seul bloc separe par \par, mais elles portent
%  deux cles : les ecrire ensemble ferait disparaitre la seconde, et
%  c'est la faute que la colonne egyptienne avait faite au meme
%  tableau.
%
%  LE MARATHI REPOND AUX DEUX NOTES, COMME IL AVAIT REPONDU A CELLE
%  DE LA SOLIVE. La premiere explique que « babo » vient de l'anglais
%  baby, du francais bebe, de l'italien bambino ; le marathi dit
%  बाळ, et le dit depuis toujours. La seconde aligne cinq langues
%  pour le « lambrequino », cette bande d'etoffe qui pend au-dessus
%  d'une cheminee ou d'un rideau ; le marathi dit झालर. Trois notes
%  du livret, trois mots que le marathi n'a pas eu a emprunter.
%
%  LA CUISINE EST LE PLUS BEAU MORCEAU DE LA COLONNE, parce que le
%  marathi a un mot propre pour chaque ustensile et qu'aucun n'est
%  emprunte :
%
%    चूल (1)      le foyer          विस्तव (3)   le feu
%    फुंकणी (8)   le soufflet       कंदील (11)   la lanterne
%    डाव (21)     la louche         झारा (22)    l'ecumoire
%    कढई (47)     la poele          हंडा (46)    la marmite
%    पातेली (26)  les casseroles    मडकी (43)    les pots
%    चाळण (28)    la passoire       किसणी (29)   la rape
%    नरसाळे (38)  l'entonnoir       मोरी (49)    l'evier
%    घागर (50)    la grande cruche  बंब (56)     le chaudron
%    खण (33)      le tiroir         फरशी (57)    le carrelage
%    रस्सा (24)   la sauce          कोयता (40)   le couperet
%
%  ET LA CHAMBRE FAIT DE MEME, avec des mots que le francais rend
%  par deux : लोड (25) le traversin, उशी (26) l'oreiller, दुलई (27)
%  l'edredon, पलंगपोस (32) le dessus-de-lit, चांदवा (18) le ciel de
%  lit, अंगुस्ताण le de a coudre, शमादान le bougeoir, सुरई (42) la
%  carafe, पाळणा (61) le berceau, फडताळ (5) le buffet.
%
%  CINQ INVERSIONS DE RENVOI PREVENUES, toujours la postposition :
%  « X dans Y » et « X de Y » sortent a l'envers. Le remede reste
%  celui des tableaux 3 a 5 — nommer d'abord ce que l'ido nomme
%  d'abord, puis qualifier par une seconde proposition.
% ===================================================================

%%K t06-apar-1 apar
\VUsaut{6.0mm}
\VUtitre{15.0pt}{60}{तक्ता क्रमांक ६}
\VUsaut{1.4mm}
\VUtitre{11.4pt}{0}{घराचा आतला भाग, सामानसुमान, अन्न,}
\VUsaut{0.4mm}
\VUtitre{11.4pt}{0}{प्रकाश.}
\VUsaut{2.2mm}

%%K t06-apar-2 apar
घर पुरे झाले आहे. जॉनचे काका आपल्या नव्या घरात राहायला आले आहेत, आणि त्या घराची रचना आपल्याला
माहीत आहे. आता ते ते घर कसे सजवतात हे पाहू. सुरुवात इथून करू :

%%K t06-tit-1 sub
\VUpk{10.2pt}{40}{शयनगृह.}
\VUsaut{3.0mm}

%%K t06-c1-01-1 p
१. --- आपण गैरसोयीच्या वेळी आलो आहोत, कारण मोलकरीण \VUgras{खोली} आवरण्यात गुंतली आहे.

%%K t06-c1-02-1 p
२. --- \VUgras{प्रसाधन-मेजावर} \textsuperscript{(1)} इथेतिथे अशा वस्तू पडल्या आहेत, ज्यांनी
माणूस हातपाय धुतो, केस विंचरतो — थोडक्यात, आवरतो. त्या म्हणजे \VUgras{घंगाळे}
\textsuperscript{(2)} आणि \VUgras{पाण्याचा तांब्या} \textsuperscript{(3)},
\VUgras{केसांचा ब्रश} \textsuperscript{(4)}, \VUgras{फणी} \textsuperscript{(5)}, \VUgras{साबण}
\textsuperscript{(6)} आणि \VUgras{स्पंज} \textsuperscript{(7)}, आणि शेवटी द्रव दंतमंजनाची
\VUgras{कुपी} \textsuperscript{(8)}.

%%K t06-c1-03-1 p
३. --- जमिनीवर, प्रसाधन-मेजासमोर, हे बघ मळके पाणी गोळा करायची \VUgras{बादली}
\textsuperscript{(9)}.

%%K t06-c1-04-1 p
४. --- \VUgras{बाहेरच्या खोलीत} \textsuperscript{(10)} आपल्याला काही \VUgras{खुंट्या}
\textsuperscript{(13)} आणि दोन \VUgras{छत्र्या} \textsuperscript{(11)} दिसतात —
\VUgras{छत्रीदाणीत} \textsuperscript{(12)}.

%%K t06-c1-05-1 p
५. --- दाराच्या दुसऱ्या बाजूला \VUgras{आरशाचे कपाट} \textsuperscript{(14)} आहे, ज्यात
\VUgras{कपडे} \textsuperscript{(15)} ठेवलेले असतात.

%%K t06-c1-06-1 p
६. --- \VUgras{मोलकरीण} \textsuperscript{(6)} \VUgras{पलंग} \textsuperscript{(17)} आवरते आहे,
तेवढ्यात आपण त्याचे वेगवेगळे भाग पाहून घेऊ. \VUgras{पलंगाच्या चौकटीत} \textsuperscript{(20)}
चांगली \VUgras{स्प्रिंगची गादी} \textsuperscript{(21)} आहे, आणि तिच्यावर युस्तिना लगेच
\VUgras{लोकरीची गादी} \textsuperscript{(22)} घालील. मग ती खुर्च्यांवरून दोन \VUgras{चादरी}
\textsuperscript{(23)} आणि \VUgras{पांघरूण} \textsuperscript{(24)} घेईल. नंतर पलंगाच्या
डोक्याशी \VUgras{लोड} \textsuperscript{(25)} आणि \VUgras{उशी} \textsuperscript{(26)} ठेवील,
ज्या \VUgras{दुलईसह} \textsuperscript{(27)} \VUgras{पलंगपोसावर} \textsuperscript{(32)} आहेत.
\VUgras{पडदे} \textsuperscript{(19)} आणि \VUgras{चांदवा} \textsuperscript{(18)} झाडायला ती
पिसांचा झाडू वापरते, आणि हे बघ कामाचा एक भाग पुरा झाला.

%%K t06-c1-07-1 p
७. --- मग तिला \VUgras{पलंगाशेजारचे लहान मेज} \textsuperscript{(28)} थोडे नीट लावावे लागेल,
\VUgras{गजराचे घड्याळ} \textsuperscript{(29)} भरावे लागेल, \VUgras{रात्रीचा दिवा}
\textsuperscript{(30)} तयार ठेवावा लागेल आणि शमादान साफ करण्यासाठी \VUgras{मेणबत्ती}
\textsuperscript{(31)} काढावी लागेल.

%%K t06-tit-2 sub
\VUsaut{5.0mm}
\VUpk{10.2pt}{40}{शिवणकामाची टोपली आणि शेकोटीचे साहित्य.}
\VUsaut{3.0mm}

%%K t06-c1-08-1 p
८. --- \VUgras{शिवणकामाची टोपलीही} \textsuperscript{(34)} आवरायला हवीच — ती \VUgras{पाटाजवळची}
\textsuperscript{(33)}. तिला \VUgras{भरतकामाची} \textsuperscript{(35)} घडी करावी लागेल,
\VUgras{शिवणकामाच्या मेजात} \textsuperscript{(38)} \VUgras{अंगुस्ताण}, \VUgras{दोरा}
\textsuperscript{(36)}, \VUgras{सुया} \textsuperscript{(37)} आणि \VUgras{कात्री}
\textsuperscript{(39)} ठेवावी लागेल, आणि \VUgras{किल्लीने} \textsuperscript{(40)} मेजाचे झाकण
बंद करावे लागेल.

%%K t06-c1-09-1 p
९. --- मधल्या \VUgras{एकपायी मेजावर} \textsuperscript{(41)} युस्तिना \VUgras{सुरईतले}
\textsuperscript{(42)} पाणी बदलेल. ती \VUgras{साखरदाणीत} \textsuperscript{(43)} \VUgras{साखर}
घालील, \VUgras{साखरेचा चिमटा} \textsuperscript{(44)} पुसील आणि \VUgras{कपड्यांचा ब्रश}
\textsuperscript{(46)} उचलून ठेवील.

%%K t06-c1-10-1 p
१०. --- खोलीची उजवी बाजू तिच्याकडून कमी काम मागेल, कारण \VUgras{लांब आरामखुर्ची}
\textsuperscript{(47)} आणि तिची \VUgras{उशी} \textsuperscript{(48)} जागच्या जागी आहेत, आणि
थंडी नसल्यामुळे \VUgras{शेकोटीच्या} \textsuperscript{(49)} साहित्यातले काहीच हललेले नाही.
\VUgras{फावडे} \textsuperscript{(50)}, \VUgras{चिमटा} \textsuperscript{(51)},
\VUgras{सरपणाचे घोडे} \textsuperscript{(55)}, \VUgras{राखेची कड} \textsuperscript{(56)} स्वच्छ
आहेत; \VUgras{शेकोटीचा पडदा} \textsuperscript{(54)} हलवलेला नाही आणि \VUgras{सरपणाची पेटी}
\textsuperscript{(52)} \VUgras{सरपणाने} \textsuperscript{(53)} भरलेली आहे.

%%K t06-noto-1 noto
\VUnotes{143.2mm}{%
(*) Babo = तान्हे मूल; इंग्रजी \textit{baby}, फ्रेंच \textit{bébé}, इटालियन \textit{bambino}.%
}

%%K t06-noto-2 noto
\VUnotes{143.2mm}{%
(*) Lambrequino = फ्रेंच \textit{lambrequin}, स्पॅनिश \textit{lambrequines}, पोर्तुगीज
\textit{lambrequins}, जर्मन आणि इंग्रजीतही \textit{lambrequins}; म्हणजे ज. इं. फ्रें. पो.
स्पॅ.%
}

%%K t06-c1-11-1 p
११. --- तरीही तिला \VUgras{लंबकाचे घड्याळ} \textsuperscript{(57)}, \VUgras{दीपदाने}
\textsuperscript{(58)} आणि \VUgras{खिशातल्या वस्तूंच्या ताटल्या} \textsuperscript{(59)}
झाडाव्या लागतील, आणि शेवटी \VUgras{आरसा} \textsuperscript{(60)} पुसावा लागेल.

%%K t06-c1-12-1 p
१२. --- राहिला तान्ह्या मुलाचा (\VUgras{बाळाचा} (*)) \VUgras{पाळणा} \textsuperscript{(61)}, तो
आधीच तयार आहे.

%%K t06-tit-3 sub
\VUsaut{5.0mm}
\VUpk{10.2pt}{40}{दिवाणखाना.}
\VUsaut{3.0mm}

%%K t06-c2-01-1 p
१. --- आता हा दिवाणखाना. फार सुंदर \VUgras{खोली} आहे. उजव्या कोपऱ्यातल्या शेकोटीवर एक नाजूक
\VUgras{मूर्ती} \textsuperscript{(1)} आणि दोन सुंदर \VUgras{फुलदाण्या} \textsuperscript{(3)}
मांडल्या आहेत, आणि तिच्यावर मखमली \VUgras{झालर} \textsuperscript{(2)} (*) सोडली आहे.

%%K t06-c2-02-1 p
२. --- चुलीतल्या ठिणग्यांनी गालिचा पेटू नये म्हणून चुलीसमोर तांब्याची \VUgras{ठिणग्यांची जाळी}
\textsuperscript{(4)} ठेवली आहे.

%%K t06-c2-03-1 p
३. --- जवळच एक देखणे बायकी \VUgras{लिहिण्याचे मेज} \textsuperscript{(5)} उघडे आहे. त्यावर
\VUgras{घरमालकीण} \textsuperscript{(23)} आपली पत्रे लिहिते.

%%K t06-c2-04-1 p
४. --- खिडकीला \VUgras{काचेवरचे पडदे} \textsuperscript{(6)} आहेत, \VUgras{बांधपट्ट्यांनी}
\textsuperscript{(8)} \VUgras{कड्यांना} \textsuperscript{(7)} अडकवलेले, आणि
\VUgras{वरून सोडलेले} \textsuperscript{(9)} जाड कापडी पडदेही आहेत.

%%K t06-c2-05-1 p
५. --- खिडकीच्या कोनाड्यात एका \VUgras{कुंडीत} \textsuperscript{(11)} हिरवे \VUgras{रोप}
\textsuperscript{(10)} आहे — एक देखणे ताड, ज्याची पाने जवळजवळ \VUgras{रॉकेलच्या दिव्यापर्यंत}
\textsuperscript{(12)} पोहोचली आहेत.

%%K t06-c2-06-1 p
६. --- दिव्याची \VUgras{छत्री} \textsuperscript{(13)} मारियाने नीट लावली आहे — ती लाघवी
\VUgras{तरुणी} \textsuperscript{(14)}, जी \VUgras{पियानोपाशी} \textsuperscript{(15)} आहे. अगदी
ताठ बसून, \VUgras{पाटावर} \textsuperscript{(16)}, ती संगीताचा अभ्यास करते. ती फार हुशार आणि
नीटनेटकी आहे, हे तिच्या \VUgras{संगीताच्या कपाटावरून} \textsuperscript{(17)} कळते — ते अगदी
व्यवस्थित लावलेले आहे.

%%K t06-tit-4 sub
\VUsaut{5.0mm}
\VUpk{10.2pt}{40}{खोडकर पोरगा.}
\VUsaut{3.0mm}

%%K t06-c2-07-1 p
७. --- तिचा भाऊ पौलुस \VUgras{पुस्तक} \textsuperscript{(19)} शोधतो आहे,
\VUgras{पुस्तकांच्या कपाटात} \textsuperscript{(18)}. तो \VUgras{तरुण} \textsuperscript{(20)}
आहे, चंचल आणि असह्य. फार वरचे पुस्तक गाठण्यासाठी कपाटाला लोंबकळताना, वर ठेवलेला
\VUgras{अर्धपुतळा} \textsuperscript{(21)} पाडायचा धोका तो पत्करतो आहे.

%%K t06-c2-08-1 p
८. --- असा खोडसाळपणा करायला त्याला ही संधी मिळाली, कारण त्याची आई एका मैत्रिणीशी — काही दिवस
राहायला आलेल्या \VUgras{बाईंशी} \textsuperscript{(22)} — बोलण्यात गुंतली आहे. आई
\VUgras{सोफ्यावर} \textsuperscript{(24)} बसली आहे, \VUgras{आरामखुर्चीजवळ}
\textsuperscript{(25)}, आणि तिची मैत्रीण त्या \VUgras{खुर्चीवर} \textsuperscript{(27)} आहे,
जिचे फक्त \VUgras{पाठटेकण} \textsuperscript{(26)} दिसते.

%%K t06-c2-09-1 p
९. --- भिंतीवर टांगलेल्या मोठ्या \VUgras{चौकटीत} \textsuperscript{(30)} एका नातेवाईक बाईची
\VUgras{तसबीर} \textsuperscript{(29)} आहे. मेजावर ठेवलेल्या \VUgras{अल्बममध्ये}
\textsuperscript{(32)} बाकीच्या घरच्यांची \VUgras{छायाचित्रे} \textsuperscript{(33)} जमवली
आहेत. मुलांनी तो नुकताच चाळला, कारण \VUgras{खुर्च्या} \textsuperscript{(31)} अजून मेजाभोवती
गोळा झालेल्या आहेत.

%%K t06-c2-10-1 p
१०. --- पोर्सिलेनची \VUgras{चिनी फुलदाणी} \textsuperscript{(34)}, जी
\VUgras{कागद कापायच्या सुरीजवळ} \textsuperscript{(35)} आहे, ती मारियाला तिच्या वाढदिवसाला भेट
मिळाली, आणि खोलीचा कोपरा सजवणारा \VUgras{आडपडदा} \textsuperscript{(36)} तिच्या काकांनी चीनहून
आणला.

%%K t06-c2-11-1 p
११. --- सुंदर \VUgras{चित्रांनी} \textsuperscript{(37)} — \VUgras{भिंतीवर}
\textsuperscript{(42)} टांगलेल्या — आनंदित झालेला आणि \VUgras{झुंबराने} \textsuperscript{(38)}
उजळलेला, ज्याचे \VUgras{विजेचे दिवे} \textsuperscript{(39)} संध्याकाळी आनंदाने चमकतात, तो
दिवाणखाना फार आल्हाददायक दिसतो. लाकडी फरशीवर पसरलेला मऊ \VUgras{गालिचा} \textsuperscript{(40)}
आणि इतकी सोयीची \VUgras{विजेची घंटी} \textsuperscript{(41)} त्याचा आराम पुरा करतात.

%%K t06-tit-5 sub
\VUsaut{3.6mm}
\VUpk{10.2pt}{40}{जेवणघर.}
\VUsaut{3.0mm}

%%K t06-c3-01-1 p
१. --- जेवणाची घंटा वाजली. मारिया सोडून सगळे घर मेजाभोवती जमले आहे. उत्तम चिनी मातीच्या
\VUgras{शेगडीने} \textsuperscript{(1)} — बारीक \VUgras{नळी} \textsuperscript{(2)} असलेल्या —
जेवणघर सुखद उबदार झाले आहे.

%%K t06-c3-02-1 p
२. --- या खोलीत फार सामान नाही. भिंतीलगत, \VUgras{लहान चौरंगाच्या} \textsuperscript{(4)} वर,
एक शोभेचे \VUgras{कारंजे} \textsuperscript{(3)} टांगले आहे. अगदी जवळ \VUgras{फडताळ}
\textsuperscript{(5)} आहे, ज्यावर थंड पदार्थ, गोड पदार्थांची ताटे आणि लगेच लागणार नाहीत अशी
मेजावरची भांडी ठेवतात.

%%K t06-c3-03-1 p
३. --- अशा रीतीने वरच्या फळीवर आपल्याला \VUgras{भाजलेले कोंबडे} \textsuperscript{(7)},
\VUgras{मासा} \textsuperscript{(8)} आणि एक \VUgras{तबक} \textsuperscript{(10)} \VUgras{फळांनी}
\textsuperscript{(9)} भरलेले दिसते. खाली हे बघ \VUgras{चीज} \textsuperscript{(11)},
\VUgras{पाव} \textsuperscript{(12)}, आणि \VUgras{तेलमिठाची तबकडी} \textsuperscript{(13)},
जिच्यात कोशिंबिरीला लागणारे सगळे असते : तेल, व्हिनेगर, \VUgras{मीठ} \textsuperscript{(14)} आणि
\VUgras{मिरी} \textsuperscript{(15)}. शेवटी, खालच्या फळीवर \VUgras{केक} \textsuperscript{(17)}
आहे, \VUgras{मोहरीच्या भांड्याजवळ} \textsuperscript{(16)}.

%%K t06-c3-04-1 p
४. --- जेवणघरातल्या घड्याळाला \VUgras{भिंतीवरचे घड्याळ} \textsuperscript{(20)} म्हणतात, आणि
त्याचा आकार फार वेगळा आहे. ते एका लाकडी चौकटीत बसवले आहे, जिच्यात शिवाय \VUgras{वायुभारमापक}
\textsuperscript{(19)} आणि \VUgras{तापमापक} \textsuperscript{(18)} आहे.

%%K t06-tit-6 sub
\VUsaut{4.6mm}
\VUpk{10.2pt}{40}{जेवणाची वाढणी.}
\VUsaut{3.0mm}

%%K t06-c3-05-1 p
५. --- खोलीच्या सगळ्या भिंतींना मेणलावलेल्या ओकच्या \VUgras{लाकडी फळ्या}
\textsuperscript{(21)} लावल्या आहेत. व्हरांड्यात जाणारे दार कापडी \VUgras{दाराच्या पडद्याने}
\textsuperscript{(22)} बऱ्यापैकी बंद होते. जेवणानंतर घर तिथे कॉफी प्यायला जाईल. \VUgras{कप}
\textsuperscript{(24)} आणि \VUgras{बशा} \textsuperscript{(25)} आधीच \VUgras{वाढपाच्या मेजावर}
\textsuperscript{(23)} तयार आहेत.

%%K t06-c3-06-1 p
६. --- मेजाच्या वर छताला \VUgras{टांगता दिवा} \textsuperscript{(26)} बसवला आहे, ज्याचा फार
तेजस्वी दिवा संध्याकाळी सगळे जेवणघर उजळून टाकतो.

%%K t06-c3-07-1 p
७. --- आता \VUgras{जेवणाचे मेजच} \textsuperscript{(27)} पाहू. घर आत आले तेव्हा वाढणी तयार
होती. \VUgras{मेजपोसावर} \textsuperscript{(28)} प्रत्येक जेवणाऱ्याच्या जागी \VUgras{ताट}
\textsuperscript{(33)}, \VUgras{रुमाल} \textsuperscript{(29)}, \VUgras{चमचा}
\textsuperscript{(34)}, \VUgras{काटा} \textsuperscript{(35)}, \VUgras{सुरी}
\textsuperscript{(36)} आणि \VUgras{पेला} \textsuperscript{(32)} मांडला होता. शिवाय दारूसाठी
\VUgras{बाटल्या} \textsuperscript{(30)} आणि पाण्यासाठी \VUgras{सुरई} \textsuperscript{(31)}
होती.

%%K t06-c3-08-1 p
८. --- \VUgras{नोकराने} \textsuperscript{(39)} आधी \VUgras{सूपचे भांडे} \textsuperscript{(37)}
आणले, ज्यात अजून थोडे सूप वाफाळते आहे. आता तो पहिला \VUgras{पदार्थ} \textsuperscript{(38)}
वाढतो आहे — मांसाचा पदार्थ. नंतर तो आपण नुकतेच नाव घेतलेले पदार्थ मांडील, आणि शेवटी,
\VUgras{गोड पदार्थानंतर}, मालक व्हरांड्यात गेल्याची संधी साधून मेजावरची भांडी उचलील आणि जेवणघर
पुन्हा आवरील.

%%K t06-c3-08-2 p
\textit{टीप.} ---
\textit{अन्नाशी संबंधित रोजच्या वापरातले शब्द इथे फार थोडक्यात दिले आहेत. ती यादी « उपाहारगृह » (क्रमांक १३) आणि « बाजार » (क्रमांक १५) या दोन तक्त्यांवर पुरी केली जाईल.}

%%K t06-tit-7 sub
\VUsaut{4.6mm}
\VUpk{10.2pt}{40}{स्वयंपाकघर.}
\VUsaut{3.0mm}

%%K t06-c4-01-1 p
१. --- अर्ध्या उघड्या दारातून आपण ज्या स्वयंपाकघरात डोकावतो, तिथे एक चित्रमय पसारा दिसतो.
\VUgras{चुलीसमोर} \textsuperscript{(1)}, जिच्यात \VUgras{विस्तव} \textsuperscript{(3)} तडतडतो
आहे, \VUgras{शीग फिरवायचे यंत्र} \textsuperscript{(5)} लावले आहे. छान \VUgras{भाजलेले मांस}
\textsuperscript{(7)} \VUgras{शिगेला लावले} \textsuperscript{(6)} आहे आणि शिजते आहे.

%%K t06-c4-02-1 p
२. --- नुकतीच वापरलेली \VUgras{फुंकणी} \textsuperscript{(8)} आणि \VUgras{कोळशाचे फावडे}
\textsuperscript{(9)} \VUgras{सरपणासारखीच} \textsuperscript{(10)} जमिनीवरच राहिली आहेत.

%%K t06-c4-03-1 p
३. --- शेकोटीचा वरचा भाग आवरलेला आहे; \VUgras{कंदील} \textsuperscript{(11)}, \VUgras{काडेपेटी}
\textsuperscript{(13)} — जिच्यात \VUgras{काड्या} \textsuperscript{(12)} आहेत —,
\VUgras{शमादान} \textsuperscript{(14)} आणि \VUgras{मेणबत्तीदाणी} \textsuperscript{(15)} आपल्या
\VUgras{मेणबत्तीसह} \textsuperscript{(16)} जागच्या जागी आहेत. पण मोठी \VUgras{कोळशाची बादली}
\textsuperscript{(18)}, \VUgras{दगडी कोळशाने} \textsuperscript{(17)} भरलेली, जी
\VUgras{स्वयंपाकिणीने} \textsuperscript{(23)} नुकतीच तळघरात भरून आणली, ती स्वयंपाकघराच्या
मधोमध राहिली आहे.

%%K t06-c4-04-1 p
४. --- खरे तर, बिचाऱ्या बाईला घाई करावी लागते, म्हणजे न्याहारी वेळेवर तयार होईल. आता ती
\VUgras{स्वयंपाकाच्या शेगडीजवळ} \textsuperscript{(19)} आहे आणि \VUgras{रस्सा}
\textsuperscript{(24)} ढवळते आहे, जो फार करपू द्यायचा नाही.

%%K t06-c4-05-1 p
५. --- \VUgras{भट्टीत} \textsuperscript{(20)} एक केक भाजतो आहे, जो तिने मुलांच्या खाऊसाठी तयार
केला. स्वयंपाकाच्या शेगडीवर \VUgras{डाव} \textsuperscript{(21)} आणि \VUgras{झारा}
\textsuperscript{(22)} टांगले आहेत.

%%K t06-tit-8 sub
\VUsaut{4.0mm}
\VUpk{10.2pt}{40}{भांडीकुंडी.}
\VUsaut{3.0mm}

%%K t06-c4-06-1 p
६. --- खोलीच्या टोकाला टांगलेली \VUgras{भांडीकुंडी} \textsuperscript{(25)} फार स्वच्छ आहे.
तांब्याची सुंदर \VUgras{पातेली} \textsuperscript{(26)}, \VUgras{दुधाचे भांडे}
\textsuperscript{(27)}, \VUgras{चाळण} \textsuperscript{(28)} आणि \VUgras{किसणी}
\textsuperscript{(29)} काळजीपूर्वक घासून ठेवली आहेत.

%%K t06-c4-07-1 p
७. --- \VUgras{मिठाची डबी} \textsuperscript{(30)} भांड्यांच्या कपाटावर
\VUgras{चहाच्या किटलीजवळ} \textsuperscript{(31)} आणि \VUgras{तेलाच्या मोठ्या बाटलीजवळ}
\textsuperscript{(32)} ठेवली आहे. \VUgras{खणात} \textsuperscript{(33)} बहुधा जेवणाची भांडी आणि
स्वयंपाकघरातल्या सुऱ्या असतील.

%%K t06-c4-08-1 p
८. --- \VUgras{झाडू} \textsuperscript{(34)} आणि \VUgras{पिसांचा झाडू} \textsuperscript{(35)}
कोपऱ्यात आहेत. स्वयंपाकिणीने नुकताच \VUgras{लाकडी ठोकळा} \textsuperscript{(36)} वापरला, तो
तिने खिडकीजवळ ठेवून दिला.

%%K t06-c4-09-1 p
९. --- \VUgras{हात पुसायचे फडके} \textsuperscript{(37)} भिंतीवर टांगले आहे,
\VUgras{नरसाळ्याजवळ} \textsuperscript{(38)}. स्वयंपाकघराच्या त्या भागात \VUgras{भाजायची जाळी}
\textsuperscript{(39)}, \VUgras{कोयता} \textsuperscript{(40)} आणि \VUgras{चिरण्याची फळी}
\textsuperscript{(41)} ठेवली आहे. पुढे, कोयत्याच्या वर, एका \VUgras{फळीवर}
\textsuperscript{(42)} \VUgras{मडकी} \textsuperscript{(43)}, \VUgras{उकळण्याचे भांडे}
\textsuperscript{(44)} आपल्या \VUgras{झाकणासह} \textsuperscript{(45)} आणि \VUgras{हंडा}
\textsuperscript{(46)} आहे. \VUgras{कढई} \textsuperscript{(47)} जवळच टांगली आहे.

%%K t06-c4-10-1 p
१०. --- या कोपऱ्यात फक्त तांब्याचा \VUgras{नळ} \textsuperscript{(48)} चमकतो आहे. \VUgras{मोरी}
\textsuperscript{(49)}, जिच्यावर जाड \VUgras{घागर} \textsuperscript{(50)} ठेवली आहे, ती आपल्या
लहानशा कपाटासह बहुधा फार सोयीची असेल — त्या कपाटात \VUgras{व्हिनेगरचे पिंप}
\textsuperscript{(51)} आपल्या \VUgras{नळासह} \textsuperscript{(52)}, \VUgras{कचऱ्याची पेटी}
\textsuperscript{(53)} आणि \VUgras{खरकटी ताटे} \textsuperscript{(54)} ठेवतात.

%%K t06-tit-9 sub
\VUsaut{4.0mm}
\VUpk{10.2pt}{40}{स्वयंपाकघरात ठेवलेल्या इतर वस्तू.}
\VUsaut{3.0mm}

%%K t06-c4-11-1 p
११. --- \VUgras{घंगाळ} \textsuperscript{(55)}, ज्यात \VUgras{भाज्या} \textsuperscript{(58)}
धुतात, आणि बिडाचा \VUgras{बंब} \textsuperscript{(56)}, \VUgras{फरशीवर} \textsuperscript{(57)}
ठेवलेला — यांचीही जागा बहुधा त्या कपाटातच असेल.

%%K t06-c4-12-1 p
१२. --- स्वयंपाकिणीला शेगड्यांची खात्रीने कमतरता नाही, कारण हे बघ शिवाय, मेजाच्या कोपऱ्यावर,
\VUgras{गॅसची शेगडी} \textsuperscript{(59)}, जिच्यावर एका लहान पातेल्यात पाणी तापते आहे.
त्यातून निघणारी \VUgras{वाफ} \textsuperscript{(60)} आपल्याला सांगते की ते बहुधा उकळते आहे. हे
पाणी बहुधा कॉफीसाठी वापरले जाईल, कारण हे बघ मेजाच्या दुसऱ्या टोकाला \VUgras{कॉफीची किटली}
\textsuperscript{(64)} आणि \VUgras{कॉफी दळायची गिरणी} \textsuperscript{(63)}.

%%K t06-c4-13-1 p
१३. --- शेगडी \VUgras{गॅसच्या बर्नरला} \textsuperscript{(62)} जोडली आहे, लांब
\VUgras{रबरी नळीने} \textsuperscript{(61)}.

%%K t06-c4-14-1 p
१४. --- स्वयंपाकिणीला मदत करायला येणारी \VUgras{रोजंदार मोलकरीण} \textsuperscript{(66)}
जेवणासाठी भाज्या तयार करते. त्या निवडून धुऊन झाल्या की ती त्या \VUgras{परातीत}
\textsuperscript{(65)} ठेवील आणि शिजवण्याच्या वेळेची वाट पाहील.

%%K t06-tit-10 sub
\VUsaut{4.0mm}
{\centering\fontsize{10.2pt}{10.2pt}\selectfont\textls[40]{\scshape न्हाणीघर.} \textit{(स्नानगृह.)}\par}
\VUsaut{3.0mm}

%%K t06-c5-01-1 p
१. --- घरातल्या मुख्य खोल्या ओळखीच्या व्हायला आता एकच अशिष्टपणा बाकी आहे : न्हाणीघराची मांडणी
पाहणे. त्याला फार वेळ लागणार नाही, कारण ते मोठे नाही, आणि आपल्याला पटकन कळेल की \VUgras{टबला}
\textsuperscript{(1)} चांगला \VUgras{पाणी तापवायचा बंब} \textsuperscript{(2)} आहे,
\VUgras{साबणदाणी} \textsuperscript{(3)} आंघोळ करणाऱ्याच्या हाताला लागेल अशी आहे, आणि
\VUgras{पंचे अडकवायची दांडी} \textsuperscript{(5)} \VUgras{पंचे} \textsuperscript{(4)} सहज
वाळवील यात शंका नाही.

%%K t06-c5-02-1 p
२. --- या खोलीच्या भिंतींना \VUgras{चिनी मातीच्या फरशा} \textsuperscript{(6)} लावल्या आहेत,
ज्यामुळे \VUgras{शॉवर} \textsuperscript{(7)} बसवणे शक्य झाले — वापरताना उडणारे पाणी भिंती खराब
करील अशी भीती न बाळगता. पायांच्या आंघोळीसाठी वापरायची जस्ती \VUgras{लहान परात}
\textsuperscript{(8)} हे सामान पुरे करते.

\VUsaut{6.0mm}
\VUfilet{22mm}
